\documentclass[12pt,a4paper]{article}

% ==========================================
% PACOTES E FORMATAÇÃO
% ==========================================

\usepackage[utf8]{inputenc}
\usepackage[T1]{fontenc}
\usepackage[brazil]{babel}
\usepackage{setspace}
\usepackage{geometry}
\usepackage{booktabs}
\usepackage{array}
\usepackage{amsmath}
\usepackage{tikz}
\usetikzlibrary{arrows.meta,positioning}
\usepackage{csquotes}
\usepackage{newtxtext,newtxmath}
\usepackage[
    colorlinks=true,
    linkcolor=black,
    urlcolor=blue
]{hyperref}
\usepackage[
    backend=biber,
    style=numeric,
    sorting=none
]{biblatex}

\addbibresource{refs.bib}

\geometry{a4paper, margin=2cm}
\onehalfspacing

\begin{document}
\pagenumbering{gobble}

% ==========================================
% CAPA
% ==========================================

\begin{titlepage}
    \centering
    \normalsize
    Departamento de Ciência da Computação -- UFSJ/DCOMP\\

    \vspace*{5cm}

    {\large
        \textbf{\MakeUppercase{Mineração de dados aplicada à dinâmica de
        preços nas CEASAs}}
    }\\[0.5cm]

    {\normalsize
        Pré-processamento, Granger, Efeito Dominó e protocolo de modelagem
    }\\[2cm]

    {\normalsize
        Autores:\\[0.3cm]
        Arthur Henrique Moreira Santos\\
        Gabriel da Silva Souza\\
        Samuel de Souza Carvalho Almeida\\
    }

    \vspace{0.5cm}
    {\normalsize Ciência da Computação -- UFSJ/DCOMP}

    \vfill

    {\large São João del-Rei -- MG}\\
    {\large 2026}
\end{titlepage}

\pagenumbering{roman}
\tableofcontents
\newpage
\pagenumbering{arabic}

% ==========================================
% RESUMO
% ==========================================

\section*{Resumo}
\addcontentsline{toc}{section}{Resumo}

Este trabalho investiga a dinâmica de preços de produtos agrícolas
comercializados em diferentes Centrais de Abastecimento (CEASAs). A etapa atual
concentra-se no Efeito Dominó, entendido como a possibilidade de movimentos de
preço observados em um entreposto contribuírem para antecipar movimentos
posteriores em outro. Os dados foram coletados por um \textit{crawler} e
organizados em um banco SQLite com mais de dois milhões de cotações. O fluxo
executado até o momento compreende a preservação do banco bruto, o
pré-processamento das cotações, a construção de séries temporais comparáveis,
o teste de causalidade de Granger e a caracterização do Efeito Dominó. Na
execução encadeada documentada, 2.388.316 registros
brutos resultaram em 2.295.374 cotações no banco derivado, das quais 1.833.870
possuem preço comum em R\$/kg. O recorte temporal carregou 175.160 cotações e
produziu 252 pares candidatos entre mercados. Dos 134 pares elegíveis, todos
atenderam aos critérios de estacionariedade; foram testadas 268 direções e 62
apresentaram significância após a correção de Benjamini--Hochberg. Esses
resultados são preliminares e representam evidência de precedência temporal e
capacidade preditiva no sentido de Granger, não uma demonstração isolada de
causalidade econômica direta. Nas 62 relações significativas foram encontrados
6.149 choques na origem; 4.815 puderam ser classificados e 1.881 apresentaram
resposta no destino dentro da janela de cinco datas de segunda a sexta. O
modelo preditivo ainda não foi treinado, mas foram definidos a construção da
base supervisionada e o protocolo de validação temporal que orientarão essa
próxima etapa.

\noindent\textbf{Palavras-chave:} CEASA; mineração de dados; séries temporais;
pré-processamento; preços agrícolas; Efeito Dominó.

% ==========================================
% INTRODUÇÃO
% ==========================================

\section{Introdução}

\subsection{Contextualização}

As CEASAs concentram a comercialização atacadista de produtos agrícolas e
produzem cotações que refletem condições de oferta, demanda, sazonalidade,
logística e características regionais. Quando observados ao longo do tempo,
esses registros permitem investigar se movimentos de preço permanecem locais
ou se aparecem posteriormente em outros mercados.

O projeto considera três frentes de investigação: propagação de preços entre
mercados, assimetria entre aumentos e reduções e períodos de elevada
volatilidade. A etapa atual concentra-se na propagação de preços, denominada
Efeito Dominó. As perguntas que orientam essas frentes são apresentadas na
Seção de Objetivos.

Uma análise desse tipo exige séries comparáveis. Entretanto, a base contém
diferenças de escrita nos produtos, múltiplas formas de apresentação,
quantidades expressas em massa ou contagem, intervalos de peso, versões da
mesma cotação e valores potencialmente suspeitos. Comparar diretamente os
preços originais poderia confundir uma mudança de embalagem com uma variação
econômica real.

Diante disso, foi construído um fluxo rastreável que parte do banco produzido
pelo \textit{crawler}, gera uma base pré-processada, organiza séries temporais,
aplica o teste de Granger e caracteriza ocorrências de Efeito Dominó. O banco
bruto não é modificado e cada etapa gera uma saída derivada auditável. Até o
momento, não foi realizado o treinamento de um modelo preditivo: essa fase será
conduzida a partir da base de eventos caracterizados e do protocolo temporal
descrito neste relatório.

O uso do termo Efeito Dominó neste projeto é uma analogia à propagação
sequencial descrita em outras áreas. Na segurança de processos, por exemplo,
um evento inicial pode produzir consequências sucessivas em outras unidades
\cite{swuste2019}. Neste trabalho, a propagação investigada não envolve
acidentes: corresponde à possível transmissão temporal de informação entre os
preços observados em mercados distintos.

% ==========================================
% OBJETIVOS
% ==========================================

\subsection{Objetivos}

\subsubsection{Perguntas de pesquisa}

O projeto foi organizado a partir das seguintes perguntas:

\begin{enumerate}
    \item \textbf{Efeito Dominó:} usar os preços de outras CEASAs ajuda a
    prever aumentos de preços em um mercado específico?
    \item \textbf{Efeito Foguete e Pena:} é possível identificar e prever
    situações em que os preços sobem mais rápido do que caem entre diferentes
    CEASAs?
    \item \textbf{Agrupamento de Volatilidade:} é possível prever períodos de
    grandes oscilações nos preços de produtos agrícolas comercializados nas
    CEASAs?
\end{enumerate}

A etapa atual do trabalho concentra-se na primeira pergunta. As demais serão
desenvolvidas posteriormente.

\subsubsection{Objetivo geral}

Preparar e analisar as cotações históricas das CEASAs para investigar a
existência de influência temporal entre entrepostos no sentido de Granger,
preservando a rastreabilidade das transformações e evitando conversões sem
evidência suficiente.

\subsubsection{Objetivos específicos}

\begin{itemize}
    \item consolidar versões associadas à mesma cotação lógica;
    \item normalizar diferenças textuais sem misturar automaticamente produtos
    ou variedades distintas;
    \item interpretar unidades, embalagens, contagens e intervalos;
    \item calcular preços por quilograma, unidade-base ou item quando possível;
    \item identificar registros inadequados para a camada analítica;
    \item sinalizar e preservar possíveis \textit{outliers}, mantendo-os fora
    da camada analítica principal sem apagar os valores originais;
    \item gerar um banco derivado reproduzível e validado;
    \item construir séries diárias comparáveis por produto e entreposto;
    \item verificar a estacionariedade das séries de retornos;
    \item testar a precedência temporal entre pares de entrepostos por meio da
    causalidade de Granger;
    \item caracterizar ocorrências de Efeito Dominó nas relações significativas;
    \item construir uma base supervisionada sem atributos posteriores ao choque;
    \item comparar classificadores por validação temporal com janelas crescentes
    e reservar o período final para uma única avaliação de teste.
\end{itemize}

% ==========================================
% TRABALHOS RELACIONADOS
% ==========================================

\subsection{Trabalhos relacionados}

\textcite{swuste2019} apresentam uma revisão histórica do conceito de efeito
dominó em fábricas e agrupamentos da indústria química. Nesse contexto, o
fenômeno é caracterizado por uma sequência na qual um evento inicial produz
efeitos capazes de atingir outras unidades. Embora o domínio seja diferente,
a noção de propagação espacial e sequencial contribui para delimitar a analogia
empregada neste projeto. Aqui, os eventos analisados são movimentos de preço,
e não acidentes industriais.

O teste proposto por \textcite{granger1969} verifica se valores passados de uma
série acrescentam informação à previsão de outra série, além da informação já
contida no próprio histórico da variável prevista. Portanto, a expressão
causalidade de Granger indica precedência temporal com ganho preditivo e não,
isoladamente, um mecanismo econômico causal.

No domínio de preços agrícolas, \textcite{oliveira2023} analisaram séries
mensais dos tomates Longa Vida e Cajá comercializados em Fortaleza e Ibiapaba,
no Ceará, entre janeiro de 2011 e dezembro de 2022. Utilizando causalidade de
Granger, o estudo identificou relação bilateral para o tomate Longa Vida e
relação unidirecional de Fortaleza para Ibiapaba no caso do tomate Cajá. Esse
resultado demonstra a aplicação do teste à transmissão regional de preços de
um produto agrícola e fornece uma referência metodológica direta para a
primeira pergunta desta pesquisa.

% ==========================================
% MÉTODOS
% ==========================================

\section{Métodos}

\subsection{Caracterização da base de dados}

A base de origem é o arquivo \texttt{data/cotacoes.sqlite}, produzido pelo
processo de coleta das cotações. O banco relaciona cotações, produtos,
apresentações, entrepostos, CEASAs, coletas e proveniências. A origem das
observações também permite distinguir os registros coletados diretamente
daqueles provenientes do complemento PROHORT.

A execução encadeada documentada partiu de um retrato com 2.388.316 cotações.
Após a seleção das coletas mais recentes e a consolidação de versões, foram
obtidas 2.380.499 identidades. A exclusão de mapeamentos insuficientes resultou
em 2.295.374
cotações gravadas no banco derivado. A Tabela~\ref{tab:diagnostico} resume os
principais indicadores observados nessa execução.

\begin{table}[htbp]
    \centering
    \caption{Indicadores da execução encadeada documentada do pré-processamento.}
    \label{tab:diagnostico}
    \small
    \begin{tabular}{lr}
        \toprule
        \textbf{Indicador} & \textbf{Quantidade} \\
        \midrule
        Cotações no retrato do banco bruto & 2.388.316 \\
        Linhas associadas às últimas coletas & 2.387.646 \\
        Identidades consolidadas & 2.380.499 \\
        Cotações excluídas por mapeamento insuficiente & 85.125 \\
        Cotações gravadas no banco derivado & 2.295.374 \\
        Identidades com múltiplas versões & 6.041 \\
        Apresentações catalogadas & 1.941 \\
        Cotações com preço comum válido & 2.091.322 \\
        Cotações com preço comum por kg & 1.833.870 \\
        Cotações na camada analítica principal & 2.090.611 \\
        \textit{Outliers} desconsiderados da camada principal & 711 \\
        Regras de conversão inferida aplicadas & 13 \\
        Cotações com conversão inferida & 3.250 \\
        \bottomrule
    \end{tabular}

    \par\smallskip
    \footnotesize
    Fonte: resumo do banco pré-processado ligado aos resultados documentados.
\end{table}

Essas quantidades descrevem o retrato ligado aos bancos de Granger e Efeito
Dominó consultados. Como o crawler pode continuar acrescentando cotações, uma
nova execução pode produzir números diferentes mesmo sem alteração dos
critérios.

% ==========================================
% VISÃO GERAL DOS MÉTODOS
% ==========================================

\subsection{Visão geral do fluxo}

A Figura~\ref{fig:fluxo-metodologico} apresenta as etapas concluídas e as duas
etapas seguintes já definidas. Cada banco derivado preserva a origem e os
parâmetros de sua execução, permitindo reproduzir o caminho entre a coleta e
os resultados estatísticos.

\begin{figure}[htbp]
    \centering
    \begin{tikzpicture}[
        node distance=0.45cm,
        etapa/.style={
            draw,
            rounded corners,
            align=center,
            minimum width=9.4cm,
            minimum height=0.8cm,
            fill=black!5
        },
        fluxo/.style={-{Latex[length=2mm]}, thick}
    ]
        \node[etapa] (crawler) {\textit{Crawler} das CEASAs\\coleta das cotações};
        \node[etapa, below=of crawler] (bruto) {Banco bruto\\\texttt{data/cotacoes.sqlite}};
        \node[etapa, below=of bruto] (preprocessamento) {Pré-processamento\\\texttt{resultados/preprocessamento/}\\\texttt{cotacoes\_preprocessadas.sqlite}};
        \node[etapa, below=of preprocessamento] (series) {Construção das séries temporais\\mediana diária e retornos logarítmicos};
        \node[etapa, below=of series] (granger) {Teste de causalidade de Granger\\\texttt{resultados/granger/granger.sqlite}};
        \node[etapa, below=of granger] (domino) {Caracterização do Efeito Dominó\\\texttt{resultados/efeitos/domino/}\\\texttt{efeito\_domino.sqlite}};
        \node[etapa, below=of domino] (base) {Base supervisionada de modelagem\\atributos disponíveis até o choque};
        \node[etapa, below=of base] (modelagem) {Comparação de classificadores\\validação temporal e teste final isolado};

        \draw[fluxo] (crawler) -- (bruto);
        \draw[fluxo] (bruto) -- (preprocessamento);
        \draw[fluxo] (preprocessamento) -- (series);
        \draw[fluxo] (series) -- (granger);
        \draw[fluxo] (granger) -- (domino);
        \draw[fluxo] (domino) -- (base);
        \draw[fluxo] (base) -- (modelagem);
    \end{tikzpicture}
    \caption{Fluxo metodológico executado e etapas seguintes definidas.}
    \label{fig:fluxo-metodologico}
\end{figure}

\subsection{Ambiente e reprodutibilidade}

O pré-processamento foi implementado em Python e executado em contêiner Docker.
As dependências são gerenciadas pelo \texttt{uv}, enquanto o Docker Compose
fornece um comando único para construção da imagem e execução do processo. O
banco bruto é aberto somente para leitura, e o resultado é gravado em um novo
arquivo, \path{resultados/preprocessamento/cotacoes_preprocessadas.sqlite}.

A substituição da saída é atômica: o novo banco é construído em um arquivo
temporário e somente substitui a versão anterior após a conclusão das
validações. Assim, uma falha durante o processamento não sobrescreve uma saída
válida nem altera a fonte original.

Essa garantia é aplicada separadamente a cada etapa. O pipeline executa
pré-processamento, Granger e Efeito Dominó em sequência, mas não confirma os
três bancos em uma transação única. Se uma etapa posterior falhar, uma saída
anterior já concluída pode ter sido atualizada, enquanto os bancos antigos das
etapas seguintes permanecem preservados. A compatibilidade da cadeia é
verificada pelos hashes, identificadores de execução e parâmetros registrados
em cada artefato.

\subsection{Política de retenção dos dados do crawler}

O banco derivado não é uma cópia integral do banco do crawler. A seleção dos
campos separa atributos da observação, informações necessárias para auditoria e
estado operacional do coletor. As tabelas \texttt{backfill\_states} e
\texttt{schema\_migrations} e campos como cursor, erro, caminho e hash do
arquivo bruto e chaves internas de ocorrência permanecem somente no banco de
origem. Eles descrevem o funcionamento da coleta, e não o preço, o produto, o
mercado, a data ou a apresentação observada. Repeti-los na tabela analítica
aumentaria a redundância e poderia confundi-los com variáveis econômicas.

Por outro lado, são preservados o identificador e as datas da coleta, os
identificadores e nomes das dimensões, os textos originais do produto e da
apresentação e os campos do complemento PROHORT. Esses dados permitem localizar
a proveniência, interpretar a unidade original e distinguir informações diretas
daquelas incorporadas pelo complemento. Portanto, os metadados operacionais
não foram apagados: apenas permanecem na fonte, que continua sendo a camada de
auditoria integral.

A tabela \texttt{cotacao\_proveniencias} também não é copiada linha a linha.
Como uma cotação pode estar ligada a mais de uma coleta, uma junção direta
multiplicaria observações. Ela é usada para selecionar a coleta mais recente e
para confirmar a proveniência, sem transformar registros de rastreamento em
novas cotações.

\subsection{Tipos de decisão sobre um registro}

Neste trabalho, exclusão e filtro não são sinônimos. Foram adotadas quatro
formas de tratamento:

\begin{itemize}
    \item \textbf{preservação}: o valor original ou sua proveniência permanece
    disponível;
    \item \textbf{consolidação}: várias versões da mesma identidade passam a
    representar uma única observação;
    \item \textbf{exclusão da tabela final}: a linha não entra em
    \texttt{cotacoes\_preprocessadas}, mas permanece no banco bruto e em
    \texttt{cotacoes\_excluidas\_mapeamento};
    \item \textbf{desconsideração analítica}: a linha continua no banco
    derivado, porém não participa de uma visão ou teste cujos requisitos não
    atende.
\end{itemize}

Essa taxonomia evita interpretar ``não usado em Granger'' como ``apagado do
banco''.

\subsection{Seleção e consolidação das cotações}

A identidade lógica é representada por \texttt{chave\_identidade}, fornecida
pelo banco do crawler. O projeto de mineração trata essa chave como contrato de
entrada e não a recalcula. Uma mesma identidade pode aparecer em diferentes
coletas ou possuir versões simultâneas com preços distintos.

Para cada identidade, o processo seleciona a coleta mais recente registrada na
tabela de proveniências. A ordenação usa a data de download, a data de
processamento e, como desempate, o maior identificador da coleta. Quando a
linha não está associada à coleta escolhida para sua identidade, ela permanece
somente no histórico do banco bruto e não segue para a consolidação. Quando a
coleta escolhida contém mais de uma versão, os preços mínimo, comum e máximo
são consolidados separadamente pela mediana. A situação de mercado somente é
mantida quando as versões concordam; em caso de conflito, fica nula. Essa regra
evita escolher arbitrariamente uma versão e reduz a influência de um valor
isolado. O banco derivado mantém uma linha por identidade e registra a
quantidade de versões e a existência de conflito.

\subsection{Normalização dos produtos}

Os nomes dos produtos recebem uma chave textual sem diferenças de caixa,
acentuação, espaços ou pontuação. A descrição original e o identificador do
produto continuam preservados. A normalização permite reconhecer grafias
equivalentes, mas não junta automaticamente variedades ou classificações que
possam apresentar comportamentos de preço diferentes.

A regra remove espaços nas extremidades, ignora caixa, decompõe caracteres
Unicode, retira acentos, substitui sequências não alfanuméricas por espaços e
reduz espaços repetidos. O nome normalizado do banco tem prioridade; o alias
original é usado apenas quando a primeira opção está vazia. Termos de variedade
não são removidos, de forma que, por exemplo, pimentão e pimentão verde
permanecem chaves distintas.

Descrições sem conteúdo utilizável, como \texttt{*}, recebem o marcador
\texttt{produto\_invalido}. Esses registros permanecem identificáveis no banco,
mas não entram na visão analítica principal.

\subsection{Tratamento das apresentações e unidades}

As apresentações são classificadas segundo sua dimensão: massa, volume,
contagem, unidade nativa ou combinação dessas formas. A apresentação original
é preservada juntamente com a regra, a fórmula, a evidência e o nível de
confiança utilizados na transformação.

Quando o peso total está explícito, a conversão é classificada como
\texttt{exata}. Quando existem quantidade de itens e peso individual, o peso
total é calculado pela multiplicação desses valores e a conversão é
\texttt{calculada}. Unidades como dúzia e cento são convertidas para 12 e 100
itens, respectivamente, mas não recebem peso em quilogramas sem evidência do
peso individual.

Volumes são preservados em litros, mas não convertidos para quilogramas porque
a fonte não informa a densidade. Unidades nativas também são mantidas em sua
própria base. O critério geral é somente criar R\$/kg quando o denominador pode
ser obtido de forma determinística ou por uma regra inferida que atenda a todos
os controles descritos a seguir.

Intervalos como \texttt{CX 25--28 KG} preservam os limites mínimo e máximo. Um
valor central não é criado, pois isso introduziria uma precisão inexistente na
fonte. O mesmo princípio é aplicado às apresentações ambíguas: quando não é
possível determinar o denominador do preço, a conversão para R\$/kg não é
realizada.

Também foi tratada uma inconsistência textual específica na qual o sufixo de
origem \texttt{PE} aparece concatenado à unidade de massa. Dessa forma,
\texttt{Cx.20 KgPE} é interpretado como 20 kg com o sufixo PE, enquanto
\texttt{Cx.23 a27KgPE} permanece como intervalo de 23 a 27 kg.

O parser também reconhece intervalos com hífen antes da unidade, como
\texttt{Cx.17a18-Kg}, preserva pesos totais explícitos quando o texto contém
um intervalo independente de itens e identifica contagens explícitas de
unidades ou vasos mesmo quando não existe unidade estruturada na fonte.

\subsection{Conversões inferidas}

Apresentações sem peso explícito podem ser comparadas com apresentações que
possuem peso declarado para a mesma chave textual de produto, assinatura da
apresentação e origem dos dados. A assinatura normaliza o texto e remove apenas
termos de massa. Uma regra somente atende aos critérios estruturais quando
candidato e âncora possuem ao menos cinco identidades e três datas, quando
todos os pesos explícitos observados concordam e quando a data limite do treino
foi informada.

Os preços não são utilizados para escolher o peso inferido. Essa decisão evita
aproximar artificialmente os preços entre CEASAs e produzir o padrão de
propagação que a pesquisa pretende investigar. Para impedir que informações
futuras participem da criação das regras, a data limite de inferência foi
fixada em 31 de outubro de 2025, a mesma data de corte do conjunto de treino.
Somente registros com entreposto definido e preço comum positivo contam como
suporte, mas o valor desse preço não participa da escolha do peso.
Na execução documentada, 13 regras foram aplicadas a 3.250 cotações; outras 11 regras
permaneceram apenas como candidatas por não atenderem a todos os critérios
definidos.

\subsection{Exclusões e indicadores de qualidade}

Na versão atual do código, cotações que permanecem com os status
\path{sem_unidade}, \path{sem_quantidade_exata} ou \path{ambigua} não são
gravadas na tabela final. As quantidades excluídas por cada motivo são
registradas em \texttt{resumo\_preprocessamento}. As próprias linhas ficam
disponíveis em \texttt{cotacoes\_excluidas\_mapeamento}, além de continuarem
preservadas no banco bruto.

Na execução documentada, as 85.125 exclusões foram compostas por 85.108
registros sem unidade, 16 sem quantidade exata e um ambíguo. Essa
distribuição é registrada para permitir que uma futura ampliação do parser
seja avaliada pelo tipo de informação que ainda falta.

Apresentações ausentes, intervalos, volumes e unidades nativas não são
automaticamente excluídos. Elas permanecem na tabela final com marcadores e
somente ficam fora de análises que exigem uma massa total exata. Essa decisão
preserva dados que ainda podem ser úteis em comparações na unidade original e
evita tratar toda indisponibilidade de R\$/kg como erro de coleta.

Preços comuns nulos, iguais a zero ou negativos são marcados como inválidos e
ficam fora de \texttt{cotacoes\_analiticas}. Quando a relação
$\text{mínimo} \leq \text{comum} \leq \text{máximo}$ não é respeitada, o preço
comum tratado é mantido, mas os limites mínimo e máximo tratados ficam nulos.
Essa decisão evita trocar ou ordenar automaticamente valores cuja coluna
incorreta não pode ser determinada.

Os possíveis \textit{outliers} são identificados dentro de grupos com o mesmo
produto, entreposto e apresentação. Com pelo menos cinco observações, preços
por kg inferiores a 0,1 vez ou superiores a 10 vezes a mediana do grupo são
sinalizados. Preço comum igual a R\$~0,01 ou razão inferior a 0,01 ou superior
a 20 recebe a classificação \texttt{erro\_provavel}. Entre os demais casos,
grupos com pelo menos cinco ocorrências suspeitas recebem
\texttt{inconsistencia\_recorrente}; os restantes são classificados como
\texttt{extremo\_isolado}.

Entre os 711 casos sinalizados na execução documentada, 175 foram classificados
como erro provável, 441 como inconsistência recorrente e 95 como extremo
isolado.

Os valores sinalizados não são apagados nem corrigidos. Eles permanecem na
tabela principal e nas visões de auditoria, mas ficam fora de
\texttt{cotacoes\_analiticas\_principais}. Assim, não existe uma etapa de
revisão manual obrigatória para que a base seja utilizada.

A etapa de Granger desconsidera todas as três classes de suspeita, e não apenas
\texttt{erro\_provavel}. O motivo é que não houve revisão externa capaz de
separar um extremo econômico real de uma falha de denominador. A cotação
continua disponível para que essa decisão possa ser revista futuramente.

A Tabela~\ref{tab:destino-registros} resume a diferença entre exclusão e
desconsideração analítica.

\begin{table}[htbp]
    \centering
    \caption{Destino dos registros conforme a regra de qualidade.}
    \label{tab:destino-registros}
    \scriptsize
    \begin{tabular}{p{3.2cm}p{4.0cm}p{6.2cm}}
        \toprule
        \textbf{Condição} & \textbf{Destino} & \textbf{Critério e motivo} \\
        \midrule
        Sem unidade, sem quantidade exata ou ambígua &
        Banco bruto e tabela de exclusões &
        Não existe denominador identificável com segurança; a linha não entra
        na tabela final, mas permanece auditável. \\
        \addlinespace
        Sem apresentação, intervalo, volume ou unidade nativa &
        Tabela final com marcador &
        O dado pode ser útil em sua unidade original; somente fica fora de
        análises que exigem massa total exata. \\
        \addlinespace
        Produto, preço comum ou dimensão inválida &
        Tabela final, fora da visão analítica &
        Falta uma dimensão essencial para interpretar a observação. \\
        \addlinespace
        \textit{Outlier} suspeito &
        Tabela final e visão de auditoria &
        O valor não é corrigido, mas não participa de Granger sem revisão
        externa. \\
        \addlinespace
        Conversão incerta &
        Tabela final, fora de Granger &
        A comparação entre mercados exige preço em uma unidade comum. \\
        \bottomrule
    \end{tabular}
\end{table}

\subsection{Estrutura e validação da saída}

O banco derivado contém a tabela principal
\path{cotacoes_preprocessadas}, as visões \path{cotacoes_analiticas} e
\path{cotacoes_analiticas_principais}, a tabela auditável de exclusões, o
catálogo de apresentações, as regras de conversão, as classificações de
\textit{outliers}, o resumo e os parâmetros da execução.

Ao final, são verificadas a integridade interna do SQLite, a unicidade das
identidades, a consistência dos marcadores, a cobertura do catálogo e a
igualdade entre as quantidades consolidadas, gravadas e excluídas. A validação
também confere a cobertura auditável das exclusões e garante que nenhum
\textit{outlier} permaneça na camada analítica principal. Essas verificações
foram concluídas para o artefato de pré-processamento ligado aos resultados
apresentados.

% ==========================================
% SÉRIES TEMPORAIS
% ==========================================

\subsection{Construção das séries temporais}

\subsubsection{Recorte e critérios de entrada}

O recorte atual considera 12 produtos e sete entrepostos: Recife, Grande
Vitória, Curitiba, Maringá, Cascavel, Foz do Iguaçu e Londrina. Foram utilizadas
cotações entre 4 de janeiro de 2022 e 31 de outubro de 2026, desde que possuam
preço comum válido em R\$/kg, não estejam sinalizadas como \textit{outlier} e
tenham conversão exata, calculada, padronizada ou inferida com alta confiança.
O período até 31 de outubro de 2025 foi destinado à construção e ao diagnóstico
das séries; o intervalo posterior foi preservado para avaliações futuras.

Antes dessa seleção, o programa exige que o banco pré-processado tenha uma
execução concluída e que o limite temporal usado para inferir pesos seja igual
ao fim do treino. Quando a origem bruta ainda está disponível, ela também não
pode ser mais recente que o banco derivado. A etapa é interrompida se um dos
sete mercados estiver ausente ou associado a mais de uma CEASA ou UF. Essas
regras impedem a combinação de dados novos com conversões antigas, o uso de
informação futura e a junção de mercados ambíguos.

Os sete entrepostos foram definidos na exploração anterior aos testes porque
apresentavam mais de mil datas, cobertura multi-anual atual e sobreposição a
partir de janeiro de 2022. Os outros 15 entrepostos do banco não foram removidos
da base derivada: ficaram fora deste recorte por terem menos de mil datas,
série encerrada antes do período ou coleta recente e curta.

Os 12 produtos formam uma cesta de interesse fixada na exploração inicial:
batata, tomate, cebola, cenoura, morango, beterraba, melão, pimentão verde,
pimentão, batata-doce, banana-prata e quiabo. A lista não foi determinada pelos
p-valores e não corresponde a um ranking automático. A cobertura efetiva de
cada combinação produto--mercados é verificada posteriormente, o que evita que
a seleção substantiva substitua o critério amostral.

\subsubsection{Agregação, preenchimento e retornos}

Quando existe mais de uma cotação para o mesmo produto, entreposto e data,
utiliza-se a mediana diária do preço comum em R\$/kg. Sábados e domingos são
desconsiderados nesta construção analítica, mas continuam preservados no banco
pré-processado. Essa separação ocorre porque uma cotação de fim de semana não é
necessariamente inválida: a exclusão atende à necessidade específica de formar
séries regulares e defasagens comparáveis para Granger e Efeito Dominó, não a
uma regra de qualidade aplicável a todas as análises. Falhas curtas podem ser
preenchidas pelo último valor observado por, no máximo, três datas consecutivas
de segunda a sexta. Lacunas maiores permanecem abertas para não criar
continuidade artificial. Em seguida, os preços são transformados em retornos
logarítmicos diários:

\begin{equation}
    r_t = \ln(P_t) - \ln(P_{t-1}),
\end{equation}

\noindent em que $P_t$ representa a mediana diária do preço. Essa transformação
direciona a análise às variações relativas, reduzindo o efeito dos diferentes
níveis de preço entre mercados.

São materializadas as 84 combinações entre 12 produtos e sete mercados,
inclusive as combinações sem dados. Cada ponto recebe um marcador que distingue
cotação realmente observada de valor preenchido. Depois da terceira data
consecutiva de segunda a sexta sem observação, inicia-se uma lacuna e a
sequência seguinte recebe outro segmento. O retorno somente é calculado entre
preços disponíveis no mesmo segmento; portanto, uma lacuna não pode produzir
um retorno ou uma defasagem artificial. O calendário empregado representa
todas as datas de segunda a sexta e não aplica uma lista de feriados nacionais,
estaduais ou municipais.
Assim, uma data de feriado sem cotação pode ser tratada como ausência e
preenchida dentro do limite de três datas, enquanto uma cotação efetivamente
registrada nessa data é mantida. Os valores preenchidos auxiliam o alinhamento
dos calendários, mas não contam para admitir um par nem podem iniciar ou
confirmar um evento de Efeito Dominó.

Na execução documentada, o recorte carregou 175.160 cotações, desconsiderou 320
registros de fim de semana e produziu 63.163 pontos diários observados e 8.545
pontos preenchidos. Nenhuma cotação inferida entrou nesse recorte específico.

% ==========================================
% GRANGER
% ==========================================

\subsection{Teste de causalidade de Granger}

\subsubsection{Formação e qualificação dos pares}

Para cada produto, foram formados todos os pares possíveis entre os sete
entrepostos, totalizando 252 pares candidatos. Um par somente é elegível quando
as duas séries possuem ao menos 400 datas observadas de segunda a sexta em
comum no período de treino. Pontos preenchidos não contam para esse limite. Pares sem observação
recebem \texttt{serie\_ausente}; os demais pares abaixo do limite recebem
\texttt{dias\_observados\_insuficientes}. Todos permanecem registrados com o
motivo. O limite corresponde a aproximadamente um ano e meio de observações em
datas de segunda a sexta e evita estimar regressões com uma sobreposição curta.
Além dos dias observados,
registra-se a sobreposição disponível com
preenchimento, mas ela não decide a elegibilidade. Nos pares aprovados, os
retornos são recalculados em segmentos comuns, que são interrompidos por uma
lacuna em qualquer uma das séries. Os retornos são então avaliados pelos testes
ADF e KPSS; o par segue para Granger apenas quando as duas séries atendem
simultaneamente aos critérios de estacionariedade.

\subsubsection{Diagnóstico de estacionariedade}

Os diagnósticos usam o maior segmento contínuo do par. Cada série deve conter
ao menos 30 retornos, não ser constante, apresentar ADF com $p<0{,}05$ e KPSS
com $p>0{,}05$. Exigir concordância dos dois testes reduz a dependência de uma
única hipótese nula. O ADF usa constante e escolhe internamente suas defasagens
pelo BIC; o KPSS usa constante e seleção automática de defasagens. Amostra
insuficiente, série constante ou erro numérico reprova o par e mantém o motivo
no banco. As duas direções recebem \texttt{nao\_estacionario} e não participam
da correção de múltiplos testes.

\subsubsection{Especificação, seleção de lags e significância}

As duas direções de cada par são testadas separadamente. São avaliadas de uma a
dez defasagens em uma amostra fixa, e a quantidade de defasagens é escolhida
pelo Critério de Informação Bayesiano (BIC). O teste F verifica se os valores
passados da série de origem acrescentam informação à previsão da série de
destino. Como várias hipóteses são avaliadas, os valores $p$ recebem correção
global de Benjamini--Hochberg, com taxa de falsas descobertas de 5\%
\cite{benjamini1995}.

Para uma direção $A \rightarrow B$ e um lag $p$, o modelo restrito contém
apenas o passado do destino:

\begin{equation}
    r^{B}_{t} = c + \sum_{i=1}^{p}\phi_i r^{B}_{t-i} + \varepsilon_t,
\end{equation}

\noindent enquanto o modelo irrestrito acrescenta o passado da origem:

\begin{equation}
    r^{B}_{t} = c + \sum_{i=1}^{p}\phi_i r^{B}_{t-i}
    + \sum_{i=1}^{p}\beta_i r^{A}_{t-i} + \varepsilon_t.
\end{equation}

O teste F compara esses modelos. Matriz sem posto completo, amostra insuficiente
ou resíduo nulo produz \texttt{erro\_modelo}, e não um valor $p$ artificial.

A amostra de todos os lags começa depois da décima defasagem, mesmo quando o
modelo avaliado usa menos lags. Isso impede que o BIC favoreça um modelo apenas
por ter recebido mais observações. Segmentos separados por lacunas não são
concatenados e, em caso de empate no BIC, escolhe-se o menor lag. Todos os dez
resultados são preservados, mas somente o valor $p$ do lag escolhido em cada
direção com status \texttt{testado} participa da correção global. A aprovação
principal exige $q<0{,}05$; significância apenas no valor $p$ bruto não envia a
relação para o Efeito Dominó.

Um resultado significativo é interpretado como evidência de influência
temporal preditiva no sentido de Granger. Ele não prova, sozinho, que um
entreposto provoque diretamente a alteração de preço no outro, pois fatores
externos ou variáveis não observadas também podem explicar a relação.

\subsubsection{Validação e rastreabilidade}

O banco de Granger registra o caminho, o tamanho e o hash SHA-256 do banco
pré-processado, o identificador da execução de origem e todos os parâmetros do
recorte. Também preserva as séries, os pares rejeitados e seus motivos, os
diagnósticos de estacionariedade, os resultados de todos os lags e os
coeficientes do modelo escolhido.

Antes da substituição do arquivo anterior, são verificadas a integridade do
SQLite e suas chaves estrangeiras, a aplicação do limite mínimo de cobertura,
a presença das duas direções para cada par elegível, os dez resultados de lag
para cada direção testada, a seleção de um lag entre 1 e 10 e a permanência dos
valores $q$ no intervalo de zero a um. Essas verificações confirmam a coerência
interna do artefato, mas não transformam associação temporal em causalidade
econômica.

\subsection{Caracterização do Efeito Dominó}

\subsubsection{Relações analisadas e definição de choque}

O resultado de Granger é utilizado como um mapa das relações que devem ser
observadas, e não como um modelo de previsão. Para cada produto, somente as
direções que permaneceram significativas após a correção FDR são examinadas.
Em seguida, as séries são percorridas cronologicamente para localizar uma
alta no mercado de origem e verificar se ela foi seguida por uma alta no
mercado de destino.

Uma alta é considerada forte quando supera simultaneamente 5\% e o
percentil 90 dos retornos daquele produto e entreposto. O percentil é
calculado apenas com o período de treino, impedindo que dados do período de
teste influenciem a regra. Tanto o movimento inicial quanto a resposta são
determinados por cotações realmente observadas; valores preenchidos para
cobrir falhas curtas não podem iniciar nem confirmar um evento.

Como os retornos são logarítmicos, o piso de 5\% é representado por
$\ln(1+0{,}05)$. Para cada produto e mercado, o limite usado é

\begin{equation}
    L = \max\left\{\ln(1+0{,}05),\;Q_{0{,}90}(r_{\mathrm{treino,obs}})\right\}.
\end{equation}

O limite é próprio de cada produto e mercado, em vez de um valor único para
toda a base, porque a distribuição das variações difere entre séries. O piso de
5\% evita classificar como choque uma variação pequena em uma série estável; o
percentil 90 adapta a exigência quando a volatilidade histórica é maior.

\subsubsection{Classificação dos eventos}

Para uma relação direcionada $A \rightarrow B$, a caracterização segue os
passos abaixo:

\begin{enumerate}
    \setlength{\itemsep}{0.2em}
    \item localizar uma alta forte do produto no entreposto $A$;
    \item confirmar que existe uma cotação observada em $B$ na mesma data;
    \item desconsiderar como propagação os casos em que $A$ e $B$ já apresentam
    uma alta forte simultaneamente;
    \item observar as cotações de $B$ nas cinco datas de segunda a sexta seguintes, sem
    atravessar a divisão entre treino e teste;
    \item registrar Efeito Dominó quando a alta acumulada em $B$ atingir o
    limite calculado para sua própria série, usando como base o preço observado
    em $B$ na data inicial.
\end{enumerate}

Quando existe ao menos uma cotação futura observada no destino, mas a resposta
não ocorre dentro da janela, o caso é marcado como não propagado. Casos sem
cotação inicial no destino, sem nenhuma observação posterior ou com movimento
simultâneo são preservados para auditoria, mas não recebem rótulo positivo ou
negativo. Essa distinção evita transformar ausência de dados em ausência de
propagação. A resposta registrada é a primeira data observada em que a alta
acumulada atinge o limite, e não necessariamente a data da maior alta na
janela.

A Tabela~\ref{tab:regra-domino} explicita a decisão atribuída a cada condição.
A janela pode conter menos de cinco cotações: uma única observação futura já
permite o rótulo negativo se o limite não for atingido; somente a ausência
total de observações futuras representa falta de cobertura.

\begin{table}[htbp]
    \centering
    \caption{Regras de classificação dos candidatos ao Efeito Dominó.}
    \label{tab:regra-domino}
    \scriptsize
    \begin{tabular}{p{4.0cm}p{2.2cm}p{7.0cm}}
        \toprule
        \textbf{Condição} & \textbf{Rótulo} & \textbf{Decisão} \\
        \midrule
        Cotação futura observada atinge o limite & 1 &
        \texttt{confirmado}; registra a primeira resposta. \\
        Cotação futura existe, mas não atinge o limite & 0 &
        \texttt{nao\_confirmado}; houve cobertura para testar a regra. \\
        Destino também supera seu limite na data inicial & Nulo &
        \texttt{movimento\_simultaneo}; não há propagação posterior. \\
        Destino sem base observada ou com preço-base inválido & Nulo &
        Não existe referência comparável para calcular a alta acumulada. \\
        Nenhuma cotação futura observada na janela & Nulo &
        \texttt{sem\_cobertura\_destino}; falta de dado não equivale a ausência de efeito. \\
        \bottomrule
    \end{tabular}
\end{table}

A unidade resultante representa um produto, uma origem, um destino e a data
da alta inicial. Assim, a nova base indicará explicitamente os casos em que a
propagação foi observada e aqueles em que ela não ocorreu. Essa adaptação segue
a ideia de um evento inicial que produz uma consequência posterior em outra
unidade, presente na literatura de efeito dominó \cite{swuste2019}, aplicada
aqui à transmissão temporal de preços.

\subsubsection{Rastreabilidade e validação}

A caracterização foi gravada em
\path{resultados/efeitos/domino/efeito_domino.sqlite}, sem modificar os bancos
anteriores. A execução pode ser feita isoladamente ou pelo pipeline completo,
que executa pré-processamento, Granger, caracterização e construção da base de
modelagem em sequência.

O banco registra o caminho e o hash SHA-256 do banco de Granger, o identificador
da execução de origem, as relações recebidas e os parâmetros usados. Antes de
substituir a saída anterior, o processo verifica a integridade do SQLite, a
correspondência entre o resumo e as quantidades de choques, eventos
classificáveis, confirmações e não confirmações, além de exigir que todo rótulo
1 tenha status \texttt{confirmado} e todo rótulo 0 tenha status
\texttt{nao\_confirmado}.

\subsection{Construção da base de modelagem}

A nova etapa foi implementada para produzir o quarto banco derivado,
\path{resultados/modelagem/base_modelagem.sqlite}, destinado ao treinamento e à
avaliação dos classificadores. O script recebe os bancos pré-processado, de
Granger e do Efeito Dominó. Antes de construir qualquer linha, verifica se as
três execuções estão concluídas, calcula seus hashes SHA-256 e confere a cadeia
de procedência: o hash do pré-processamento deve ser o registrado pelo Granger,
e o hash do Granger deve ser o registrado pelo Efeito Dominó. Os identificadores
das execuções também devem coincidir.

O banco pré-processado participa como origem indireta e como elemento de
rastreabilidade. Suas cotações individuais não são juntadas novamente aos
eventos, pois várias apresentações podem existir para o mesmo produto, mercado
e dia, o que duplicaria linhas e voltaria a misturar preços não comparáveis. Os
atributos diários usam as séries já agregadas pelo Granger, cuja mediana,
conversões e contagens foram produzidas a partir desse banco pré-processado.

\subsubsection{Unidade de análise, alvo e elegibilidade}

Cada linha representa um evento candidato identificado por produto, mercado de
origem, mercado de destino e data do choque. A tabela do Efeito Dominó fornece
as linhas e a variável-alvo:

\begin{itemize}
    \item \texttt{efeito\_domino}=1 quando existe resposta confirmada no destino
    dentro da janela de cinco datas de segunda a sexta;
    \item \texttt{efeito\_domino}=0 quando existe base e cobertura posterior
    observadas, mas a resposta não alcança o limite;
    \item rótulo nulo quando o evento é simultâneo, não existe base comparável
    no destino ou não existe observação futura que permita classificá-lo.
\end{itemize}

Somente os rótulos 0 e 1 recebem
\texttt{elegivel\_modelagem}=1 e podem entrar nas partições. Os eventos nulos
permanecem no banco para auditoria, sem serem convertidos artificialmente em
exemplos negativos. O campo \texttt{grupo\_choque}, formado por produto, origem
e data do choque, identifica linhas que nasceram do mesmo movimento inicial,
mas não as colapsa: cada destino continua representando uma pergunta preditiva
distinta.

\subsubsection{Construção dos atributos}

Os atributos são calculados separadamente para origem e destino e usam somente
informações disponíveis até a data do choque. O script produz 81 preditores,
organizados na Tabela~\ref{tab:atributos-modelagem}.

\begin{table}[htbp]
    \centering
    \caption{Grupos de atributos construídos para a base de modelagem.}
    \label{tab:atributos-modelagem}
    \small
    \begin{tabular}{p{3.4cm}p{5.3cm}p{6.0cm}}
        \toprule
        \textbf{Grupo} & \textbf{Atributos} & \textbf{Regra e finalidade} \\
        \midrule
        Identificação da relação &
        Produto, origem, destino, UFs e indicador de mesmo estado &
        Caracterizam a relação sem usar o resultado futuro; produto e mercados
        serão tratados como categorias pelo classificador. \\

        Calendário &
        Ano, mês, trimestre e dia da semana &
        Derivados exclusivamente da data do choque para representar tendência
        temporal e sazonalidade. \\

        Relação de Granger &
        Lag selecionado, estatística F, p-valor e q-valor FDR &
        Atributos fixos da relação significativa. Nas validações internas,
        permanecem condicionados à rede definida no desenvolvimento completo. \\

        Choque e limites &
        Alta logarítmica e percentual, limites da origem e do destino, excesso
        logarítmico e razão entre alta e limite &
        Medem a intensidade do choque em relação ao critério histórico usado
        para defini-lo. \\

        Preços na data &
        Preços em R\$/kg, logaritmos e diferença entre os logaritmos &
        Usam somente os preços da origem e do destino observados na data do
        choque; nenhum preço posterior participa. \\

        Qualidade da observação &
        Quantidades de cotações exatas, calculadas, padronizadas e inferidas &
        Registram quanto e que tipo de evidência compôs o preço diário em cada
        mercado. \\

        Defasagens &
        Retornos anteriores de 1 a 10 datas de segunda a sexta, para origem e destino &
        Obtidos por deslocamento em ordem crescente dentro do mesmo segmento
        contínuo; o retorno do próprio dia é armazenado separadamente. \\

        Janelas históricas &
        Média, desvio-padrão e soma dos retornos em 5, 10 e 20 datas de segunda a sexta &
        Janelas encerradas na data do choque e calculadas dentro do mesmo
        segmento contínuo, com no mínimo uma observação, para não atravessar
        lacunas longas da série. \\

        Cobertura histórica &
        Proporção de dias realmente observados nas janelas de 5, 10 e 20 dias &
        Distingue histórico observado de valores preenchidos e informa a
        qualidade recente da série. \\
        \bottomrule
    \end{tabular}
\end{table}

Valores ausentes nos primeiros dias de uma série ou depois de lacunas são
preservados; não são preenchidos globalmente durante a criação do banco.
Imputação, padronização e codificação categórica serão ajustadas novamente
somente com o treino de cada divisão e então aplicadas à respectiva validação
ou teste.

\subsubsection{Separação entre preditores e auditoria}

A tabela \texttt{eventos\_modelagem} contém identificadores, preditores, alvo e
o indicador de elegibilidade. Data de resposta, dias até a resposta, maior alta
futura, quantidade de observações futuras, status do evento e demais campos
calculados após o choque são gravados somente em
\texttt{auditoria\_eventos}. Essa separação física reduz o risco de inclusão
acidental de informação futura na matriz de atributos.

O banco também contém \texttt{catalogo\_atributos}, que registra para cada
coluna seu papel, tipo, origem, regra de construção e se pode entrar no modelo;
\texttt{fontes\_dados}, com caminho, hash, execução e versão de cada entrada;
\texttt{particoes\_modelagem}, com o papel temporal de cada evento; e
\texttt{resumo\_modelagem}, com as contagens da execução. O arquivo é construído
em um SQLite temporário, validado por integridade e coerência das contagens e
somente então substitui a saída anterior.

\subsubsection{Protocolo de validação temporal}

Será adotada validação temporal com janelas crescentes, seguindo o mesmo desenho
geral empregado no trabalho anterior de previsão do top 10 da Fórmula 1. A
primeira divisão treinará com eventos de 2022 e 2023 e validará em 2024. A
segunda treinará com eventos de 2022 a 2024 e validará de janeiro a outubro de
2025. Não haverá distribuição aleatória dos registros.

Como o rótulo do Efeito Dominó observa as cinco datas de segunda a sexta
posteriores ao choque, cada evento recebe a data nominal em que seu rótulo
estaria totalmente disponível. A regra de purga exclui de cada divisão todo evento cuja janela
ultrapasse o fim do período ao qual foi atribuído. Assim:

\begin{itemize}
    \item eventos de treino que alcançam a validação são excluídos daquele
    treino;
    \item eventos de validação cuja janela termina depois do fim da validação
    também são excluídos;
    \item no ajuste final, eventos anteriores a novembro de 2025 cuja janela
    atravessa o corte não entram no treino;
    \item eventos do teste cuja janela ultrapassa 31 de outubro de 2026 não
    entram na avaliação final.
\end{itemize}

Não há distribuição aleatória. Eventos com o mesmo
\texttt{grupo\_choque} compartilham produto, origem e data e, por consequência,
permanecem no mesmo lado da fronteira temporal. O agrupamento é preservado como
identificador para análises de dependência ou ponderação posteriores; a base
atual não elimina destinos distintos associados ao mesmo choque.

São criadas views específicas para
\texttt{treino\_validacao\_1}, \texttt{validacao\_1},
\texttt{treino\_validacao\_2}, \texttt{validacao\_2},
\texttt{treino\_final} e \texttt{teste\_final}. Todas contêm apenas eventos
classificáveis que passaram pela regra temporal correspondente.

Regressão Logística, Árvore de Decisão, Random Forest, Gradient Boosting e SVM
serão comparados sob as mesmas divisões. A configuração de cada algoritmo será
escolhida pelo F1 médio das duas validações; precisão, revocação, MCC, PR-AUC e
matrizes de confusão serão registradas como medidas complementares. Depois da
seleção, a configuração escolhida será treinada novamente com todos os eventos
até 31 de outubro de 2025 e avaliada uma única vez entre 1º de novembro de 2025
e 31 de outubro de 2026. O teste final não será usado para modificar atributos,
hiperparâmetros ou limiar de classificação.

% ==========================================
% RESULTADOS
% ==========================================

\section{Resultados}

\subsection{Resultados do teste de Granger}

A Tabela~\ref{tab:granger-resultados} resume a execução documentada. Dos 252
pares candidatos, 134 atenderam ao mínimo de sobreposição e todos passaram
pelos diagnósticos de estacionariedade. Isso produziu 268 direções testadas,
das quais 80 seriam significativas sem ajuste e 62 permaneceram significativas
após a correção global.

\begin{table}[htbp]
    \centering
    \caption{Resumo dos resultados atuais do teste de Granger.}
    \label{tab:granger-resultados}
    \small
    \begin{tabular}{lr}
        \toprule
        \textbf{Indicador} & \textbf{Quantidade} \\
        \midrule
        Cotações carregadas no recorte & 175.160 \\
        Pares candidatos & 252 \\
        Pares com série ausente & 53 \\
        Pares com cobertura observada insuficiente & 65 \\
        Pares elegíveis e estacionários & 134 \\
        Direções testadas & 268 \\
        Direções significativas sem correção & 80 \\
        Direções significativas após correção FDR & 62 \\
        Pares com ao menos uma direção significativa & 50 \\
        Relações unidirecionais & 38 \\
        Relações bidirecionais & 12 \\
        \bottomrule
    \end{tabular}

    \par\smallskip
    \footnotesize
    Fonte: banco \path{resultados/granger/granger.sqlite} da execução documentada.
\end{table}

\subsection{Resultados da caracterização do Efeito Dominó}

A Tabela~\ref{tab:efeito-domino-resultados} resume a execução documentada.
Entre os 4.815 eventos classificáveis, 1.881 foram confirmados, correspondendo
a 39,1\%. A proporção foi semelhante no treino e no teste, com 39,0\% e 39,2\%,
respectivamente.

\begin{table}[htbp]
    \centering
    \caption{Resumo da caracterização atual do Efeito Dominó.}
    \label{tab:efeito-domino-resultados}
    \small
    \begin{tabular}{lr}
        \toprule
        \textbf{Indicador} & \textbf{Quantidade} \\
        \midrule
        Relações de Granger analisadas & 62 \\
        Choques encontrados na origem & 6.149 \\
        Eventos classificáveis & 4.815 \\
        Efeitos confirmados & 1.881 \\
        Efeitos não confirmados & 2.934 \\
        Movimentos simultâneos & 841 \\
        Eventos sem base observada no destino & 470 \\
        Eventos sem cobertura posterior & 23 \\
        \bottomrule
    \end{tabular}

    \par\smallskip
    \footnotesize
    Fonte: banco
    \path{resultados/efeitos/domino/efeito_domino.sqlite} da execução documentada.
\end{table}

Todas as 62 relações apresentaram ao menos um evento classificável e uma
confirmação. Entre as confirmações, 610 ocorreram na primeira data de segunda a
sexta posterior ao choque e 449 na segunda.

\subsection{Distribuição temporal para a modelagem}

A Tabela~\ref{tab:divisoes-modelagem} apresenta a quantidade de eventos
classificáveis em cada recorte antes da purga das fronteiras. Essas quantidades
descrevem a disponibilidade bruta no banco do Efeito Dominó, não o número final
de linhas que será entregue pelas views de modelagem. A taxa positiva permaneceu
entre 37,4\% e 40,8\%, sem desbalanceamento extremo entre os períodos.

\begin{table}[htbp]
    \centering
    \caption{Distribuição bruta dos eventos antes da purga temporal.}
    \label{tab:divisoes-modelagem}
    \scriptsize
    \begin{tabular}{p{5.4cm}rrrr}
        \toprule
        \textbf{Divisão} & \textbf{Total} & \textbf{Positivos} &
        \textbf{Negativos} & \textbf{Positivos (\%)} \\
        \midrule
        Treino da divisão 1: 2022--2023 & 1.651 & 673 & 978 & 40,8 \\
        Validação da divisão 1: 2024 & 1.212 & 461 & 751 & 38,0 \\
        Treino da divisão 2: 2022--2024 & 2.863 & 1.134 & 1.729 & 39,6 \\
        Validação da divisão 2: jan.--out. de 2025 & 1.018 & 381 & 637 & 37,4 \\
        Treino final: até out. de 2025 & 3.881 & 1.515 & 2.366 & 39,0 \\
        Teste final: nov. de 2025--out. de 2026 & 934 & 366 & 568 & 39,2 \\
        \bottomrule
    \end{tabular}

    \par\smallskip
    \footnotesize
    Fonte: banco
    \path{resultados/efeitos/domino/efeito_domino.sqlite}. As contagens
    definitivas deverão ser obtidas de
    \path{resultados/modelagem/base_modelagem.sqlite} após sua execução.
\end{table}

Antes da purga, as duas validações contêm eventos dos dez produtos e das 62
relações. Na primeira, 61 relações apresentam as duas classes; na segunda,
todas apresentam ao menos um evento positivo e um negativo. Entretanto, a
quantidade por relação é pequena em alguns casos: os mínimos são três eventos
em 2024 e cinco entre janeiro e outubro de 2025. O recorte bruto do teste final
contém 58 das 62 relações, portanto os resultados agregados desse período não
representarão quatro relações sem evento classificável. Essas coberturas deverão
ser recalculadas após a materialização das views.

% ==========================================
% DISCUSSÃO
% ==========================================

\section{Discussão}

\subsection{Interpretação das relações temporais}

Os resultados mostram, no recorte analisado, evidências de influência temporal
entre determinados entrepostos. Tomate, batata e morango concentraram,
respectivamente, 14, 13 e 11 direções significativas. Curitiba e Maringá foram
as origens mais frequentes entre as relações significativas, com 18 e 15
direções. A interpretação permanece restrita à precedência temporal preditiva
no sentido de Granger; fatores externos ou variáveis não observadas podem
explicar parte das relações.

\subsection{Suficiência e dependência das observações}

O recorte bruto reúne 3.881 eventos no treino final, quantidade próxima da que
restará após a purga e suficiente para planejar uma comparação inicial de
classificadores tabulares sob um modelo global. Essa suficiência deverá ser
confirmada com a contagem materializada. A estabilidade da proporção positiva e
a presença das duas classes nas validações favorecem esse desenho. O volume,
porém, não sustenta modelos independentes para cada relação, produto ou par de
mercados, nem justifica inicialmente redes neurais ou buscas muito amplas de
hiperparâmetros.

As linhas não são completamente independentes. Antes da purga, o treino reúne
2.211 choques de origem distintos distribuídos em 798 datas. Além disso, foram
identificados 433 grupos nos quais mais de um choque foi associado ao mesmo
produto, destino e data de resposta, totalizando 647 registros adicionais. A
base preserva cada destino como uma observação distinta e registra o
\texttt{grupo\_choque} para identificar eventos que compartilham produto,
origem e data. A separação temporal impede que o mesmo grupo atravesse treino e
avaliação, mas ainda será necessário verificar, durante o treinamento, se
ponderação por grupo ou análise de sensibilidade altera as métricas.

\subsection{Limitações e próxima etapa}

A decisão de preservar cotações de fim de semana no pré-processamento e
desconsiderá-las somente nas séries temporais é provisória. Depois da análise
dos demais efeitos, essa regra será reavaliada considerando as necessidades de
todas as etapas. Até essa revisão, os registros não serão excluídos da base
pré-processada, pois podem ser válidos para outras análises. Também permanece
como trabalho futuro avaliar calendários específicos por entreposto para tratar
feriados nacionais, estaduais e municipais sem remover datas em que um mercado
tenha operado normalmente.

A futura avaliação medirá a capacidade de prever a propagação entre relações
previamente selecionadas por Granger. Ela não demonstrará, por si só, que o
mesmo desempenho se estende a pares sem relação significativa. Será necessário
comparar posteriormente os resultados com um controle formado por relações não
significativas ou datas embaralhadas.

As duas validações internas avaliarão o classificador condicionado à rede de
relações e aos limites de choque definidos no período de desenvolvimento até
outubro de 2025. Como essas etapas anteriores não serão reestimadas dentro de
cada divisão, seus resultados não constituirão uma validação temporal ponta a
ponta de todo o pipeline. Para isso, seria necessário refazer a seleção das
relações, a correção FDR e os limites de choque usando somente o treino de cada
divisão. O teste final, por outro lado, permanece isolado: nenhum dado posterior
a 31 de outubro de 2025 foi usado para definir a rede, os limites, os atributos
ou o modelo que será avaliado entre novembro de 2025 e outubro de 2026.

Até esta etapa, foram concluídos o pré-processamento, a construção das séries,
o teste de Granger e a caracterização dos eventos. O script de construção da
base de modelagem e de suas partições temporais também foi implementado, mas o
banco ainda precisa ser materializado e validado com os artefatos atuais. O
modelo preditivo ainda não foi treinado. Após essa execução, as contagens
pós-purga deverão substituir as quantidades brutas apresentadas na
Tabela~\ref{tab:divisoes-modelagem}.

\clearpage
\printbibliography[heading=bibintoc,title={Referências}]

\end{document}
